\section{问题描述}
\label{sec:model}
\subsection{系统模型与通信协议}
考虑如下 $n$ 维线性高斯系统，其中 $n\ge1$ 为固定有限维数：
\begin{equation}
 {\bm{x}}_{t+1}={\bm{A}}{\bm{x}}_t+{\bm{w}}_t,\qquad {\bm{y}}_t={\bm{x}}_t+{\bm{v}}_t\in\R^n.
 \label{eq:model}
\end{equation}
其中，${\bm{x}}_t$ 为系统状态，${\bm{y}}_t$ 为全状态含噪测量，${\bm{w}}_t$ 和 ${\bm{v}}_t$ 分别为过程噪声与测量噪声。
设 ${\bm{x}}_0\sim N({\bm{m}}_0,{\bm{P}}_0)$，${\bm{w}}_t\sim N(0,{\bm{Q}})$，${\bm{v}}_t\sim N(0,{\bm{R}})$，
${\bm{A}}$ 可逆，${\bm{P}}_0,{\bm{R}},{\bm{Z}}\succ0$，${\bm{Q}}\succeq0$。初值、各时刻噪声及触发随机数相互独立。
下文 ${\bm{I}}$ 为单位矩阵，$\|\cdot\|_2$ 对向量表示欧氏范数、对矩阵表示诱导谱范数，
$\operatorname{tr}$ 表示矩阵的迹，$\lambda_{\min}$ 表示对称矩阵的最小特征值。

图\ref{fig:protocol}给出了系统的通信与攻击过程。传感器通过可信及时反馈获得接收端实际预测 ${\bm{m}}_t^-$。
反馈内容未被篡改，但其中的预测值仍可能含有先前载荷攻击造成的偏差。
传感器以私有随机数 $U_t\sim\operatorname{Unif}[0,1]$ 生成发送事件：
\begin{equation}
 \begin{aligned}
 \gamma_t=0\ &\Longleftrightarrow\ U_t\le g({\bm{y}}_t-{\bm{m}}_t^-),\\
 g({\bm{z}})&=\exp(-{\bm{z}}^\top {\bm{Z}}{\bm{z}}/2).
 \end{aligned}
 \label{eq:trigger}
\end{equation}
$\gamma_t=0$ 表示静默，$\gamma_t=1$ 表示发送。事件位可信，攻击者在发送决定之后修改载荷，接收值为
\begin{equation}
 \widetilde {\bm{y}}_t={\bm{y}}_t+{\bm{a}}_t,\qquad \|{\bm{a}}_t\|_2\le b.
 \label{eq:peak}
\end{equation}
其中，${\bm{a}}_t$ 为载荷修改量，$b$ 为逐路径的欧氏峰值上限，静默时令 ${\bm{a}}_t=0$。
本文考虑无丢包、无时延的状态估计过程，估计值不参与物理系统的反馈控制。

\begin{ecustfigure}
\centering
\begin{tikzpicture}[>=Stealth,font=\fontsize{11}{13}\selectfont,
 box/.style={draw,line width=.6pt,rounded corners=2pt,align=center,minimum height=12mm,text width=36mm,inner xsep=2pt},
 edge/.style={->,line width=.7pt},
 edge label/.style={font=\fontsize{10.5}{12.5}\selectfont}]
\node[box] (sensor) {Sensor and\\random trigger};
\node[box,right=18mm of sensor] (attack) {Payload modification\\$\widetilde {\bm{y}}_t={\bm{y}}_t+{\bm{a}}_t$};
\node[box,right=18mm of attack] (receiver) {Receiver\\$({\bm{m}}_t^+,{\bm{P}}_t^+)$};
\node[box,below=20mm of attack] (logger) {Remote logger\\$({\bm{P}}_k^+,\gamma_{k+1})$};
\draw[edge] (sensor) -- node[above,edge label]{sent ${\bm{y}}_t$} (attack);
\draw[edge] (attack) -- (receiver);
\draw[edge] (sensor.north) -- ++(0,7mm) -| node[pos=.25,above,edge label]{trusted event $\gamma_t$} (receiver.north);
\draw[edge] (receiver.south) -- ++(0,-39mm) -| node[pos=.25,below,edge label]{trusted actual prediction feedback ${\bm{m}}_{t+1}^-$} (sensor.south);
\draw[edge] (sensor.south east) -- ++(3mm,-9mm) -| node[pos=.25,above,edge label]{events} (logger.north west);
\draw[edge] (receiver.south west) -- ++(-3mm,-9mm) coordinate (covturn)
 -- node[midway,above,edge label,xshift=-2mm]{nominal ${\bm{P}}_k^+$} (logger.north east |- covturn) -- (logger.north east);
\end{tikzpicture}
\figcaption{实际预测反馈下的载荷攻击与可信记录}{Payload attacks and trusted records with feedback of the receiver's actual prediction}
\label{fig:protocol}
\end{ecustfigure}


\subsection{状态估计与攻击模型}
记 ${\bm{R}}_1={\bm{R}}$、${\bm{R}}_0={\bm{R}}+{\bm{Z}}^{-1}$。对正定矩阵 ${\bm{P}}$，定义
\begin{equation}
 \begin{aligned}
 {\bm{K}}_i({\bm{P}})&={\bm{P}}({\bm{P}}+{\bm{R}}_i)^{-1},\\
 F_i({\bm{P}})&={\bm{P}}-{\bm{P}}({\bm{P}}+{\bm{R}}_i)^{-1}{\bm{P}}\\
       &=({\bm{P}}^{-1}+{\bm{R}}_i^{-1})^{-1},\quad i\in\{0,1\}.
 \end{aligned}
 \label{eq:gains}
\end{equation}
接收端按正常模型计算
\begin{equation}
 \begin{aligned}
 {\bm{m}}_t^+&={\bm{m}}_t^-+\gamma_t{\bm{K}}_1({\bm{P}}_t^-)(\widetilde {\bm{y}}_t-{\bm{m}}_t^-),\\
 {\bm{P}}_t^+&=F_{\gamma_t}({\bm{P}}_t^-),\\
 {\bm{m}}_{t+1}^-&={\bm{A}}{\bm{m}}_t^+,\quad {\bm{P}}_{t+1}^-={\bm{A}}{\bm{P}}_t^+{\bm{A}}^\top+{\bm{Q}},
 \end{aligned}
 \label{eq:receiver}
\end{equation}
静默时按 ${\bm{m}}_t^+={\bm{m}}_t^-$ 更新，发送时使用收到的载荷。
初值为 $({\bm{m}}_0^-,{\bm{P}}_0^-)=({\bm{m}}_0,{\bm{P}}_0)$。${\bm{P}}_t^\pm$ 是接收端按事件历史递推的名义协方差，
受攻时还需计入估计偏差，才能得到真实均方误差。
对任一有限时刻，全部事件历史对应的 ${\bm{P}}_t^\pm$ 均正定，且构成有限值集。
增益矩阵 ${\bm{K}}_i$ 一般不对称，后文的范数估计不要求 $\|{\bm{K}}_i\|_2\le1$。

用 $S$ 表示独立于物理过程和触发随机数的攻击私有种子，
令 $Y_t^{\rm true}={\bm{y}}_t$ 在发送时成立，静默时为 $\bot$。定义
\begin{equation}
 \begin{aligned}
 \mathcal G_t^-&=\sigma(S,\gamma_{0:t-1},Y_{0:t-1}^{\rm true}),\\
 \mathcal G_t^+&=\mathcal G_t^-\vee\sigma(\gamma_t,Y_t^{\rm true}).
 \end{aligned}
 \label{eq:information}
\end{equation}
攻击策略为任意 $\mathcal G_t^+$ 可测且满足式\eqref{eq:peak}的随机因果策略。
给定策略后，攻击者可由上述历史重建已实施的动作及接收端预测。
攻击者不能读取隐藏状态、未发送测量、未来噪声或触发私有随机数。
式\eqref{eq:information}中的扩展历史用于分析真实状态的条件分布。
远端评估者没有真实已发送测量的副本，仅能利用下文定义的通信记录。

\subsection{攻击损伤与统计隐蔽约束}
固定目标时刻 $k\in\{0,1,\ldots\}$。下标 $0$ 表示无攻击协议的分布，$\pi$ 表示受攻协议的分布。
两种运行均根据各自的实际预测形成发送事件，事件轨迹因而可以不同。
以真实总均方误差（mean-square error，MSE）相对无攻击运行的增量衡量攻击损伤：
\begin{equation}
 J_k(\pi)=\E_\pi\|{\bm{x}}_k-{\bm{m}}_k^+\|_2^2-
             \E_0\|{\bm{x}}_k-{\bm{m}}_k^+\|_2^2.
 \label{eq:damage}
\end{equation}
令 $H_t$ 收集截至 $t$ 的可信事件、实际预测反馈及收到的载荷（静默记为 $\bot$）。
日志保存实际发生的事件、反馈和接收值，因此其中的载荷可以包含攻击修改。
为比较不同记录条件下的损伤保证，分别考虑完整日志、事件串及协方差标记记录：
\begin{equation}
 \begin{aligned}
 R_{\rm full}&=H_{k+1},\\
 R_{\rm bits}&=\gamma_{0:k+1},\\
 R_{\rm marked}&=({\bm{P}}_k^+,\gamma_{k+1}).
 \end{aligned}
 \label{eq:records}
\end{equation}
其中，${\bm{P}}_k^+$ 为未经量化的完整协方差矩阵标签。三类记录均包含目标时刻后的下一事件，
标记记录中的协方差用于确定该事件的正常概率基准。

用 $\TV(\mu,\nu)=\sup_E|\mu(E)-\nu(E)|$ 记总变差，并定义
\begin{equation}
 \begin{aligned}
 \Delta_j(\pi)&=\TV(\Law_\pi(R_j),\Law_0(R_j)),\\
 V_j(b,\varepsilon)&=\sup_{\pi:\,\|{\bm{a}}_t\|_2\le b,\,\Delta_j(\pi)\le\varepsilon}J_k(\pi).
 \end{aligned}
 \label{eq:value}
\end{equation}
式\eqref{eq:value}以峰值 $b$ 和记录联合分布的TV预算 $\varepsilon$ 同时限制攻击，
$V_j$ 表示满足这两项约束的最坏损伤。该隐蔽约束不依赖具体检测器的报警阈值。
完整日志确定事件串，事件串进一步确定精确标记记录，故数据处理性质给出
$\Delta_{\rm marked}\le\Delta_{\rm bits}\le\Delta_{\rm full}$。

\section{双预算约束下的损伤分析}
\label{sec:bounds}
载荷修改一方面引入估计偏差，另一方面通过实际预测反馈改变后续发送概率。
由于名义协方差由事件历史决定，计算真实损伤时需要同时考虑这两种影响。
为此，先给出条件误差分解，再建立峰值约束与记录约束下的损伤上界，并构造达到相同阶数的攻击。
\begin{lemma}[条件误差分解]
\label{lem:gaussian}
对任意允许策略，${\bm{x}}_t\mid\mathcal G_t^\pm$ 为高斯分布，均值记为 ${\bm{\mu}}_t^\pm$，
协方差为 ${\bm{P}}_t^\pm$。令 ${\bm{d}}_t={\bm{\mu}}_t^--{\bm{m}}_t^-$、${\bm{D}}_t={\bm{\mu}}_t^+-{\bm{m}}_t^+$，则
\begin{equation}
 \begin{aligned}
 {\bm{D}}_t&=({\bm{I}}-{\bm{K}}_{\gamma_t}({\bm{P}}_t^-)){\bm{d}}_t-\gamma_t {\bm{K}}_1({\bm{P}}_t^-){\bm{a}}_t,\\
 {\bm{d}}_{t+1}&={\bm{A}}{\bm{D}}_t,\quad {\bm{d}}_0=0,\\
 J_k(\pi)&=\E_\pi\|{\bm{D}}_k\|_2^2+
 \E_\pi\operatorname{tr}{\bm{P}}_k^+-\E_0\operatorname{tr}{\bm{P}}_k^+.
 \end{aligned}
 \label{eq:decomposition}
\end{equation}
\end{lemma}
\begin{proof}
发送时，给定真实测量后，发送事件的似然因子不再依赖状态，
故在后验归一化时消去，条件分布按标准高斯公式更新。
静默时，将触发概率对测量噪声积分，所得似然正比于
$\exp[-({\bm{x}}_t-{\bm{m}}_t^-)^\top {\bm{R}}_0^{-1}({\bm{x}}_t-{\bm{m}}_t^-)/2]$。
因此静默后的条件均值为 ${\bm{\mu}}_t^-+{\bm{K}}_0({\bm{P}}_t^-)({\bm{m}}_t^--{\bm{\mu}}_t^-)$，
协方差为 $F_0({\bm{P}}_t^-)$。将条件均值与接收端估计相减，得到偏差递推。
${\bm{x}}_t-{\bm{\mu}}_t^+$ 在给定 $\mathcal G_t^+$ 后均值为零，
误差平方的交叉项因此具有零期望，得到最后一式。
完整的归纳证明见附录\ref{app:conditional}。
\end{proof}

式\eqref{eq:decomposition}中，第一项为条件偏差引起的损伤，后两项之差来自事件历史概率的变化。
同一事件历史下的协方差递推不变，但不同历史的概率权重可以改变，因此协方差变化项可能为负。
单个策略的 $J_k$ 也可能为负；最坏值 $V_j$ 的可行策略包含零动作策略，故 $V_j\ge0$。
在此基础上，可得以下双预算损伤界。

\Needspace{13\baselineskip}
\begin{theorem}[任意固定维数下的最坏损伤界]
\label{thm:main}
固定式\eqref{eq:model}--\eqref{eq:information}的模型、有限 $k$ 及参考峰值
$b_{\rm ref}>0$。存在 $0<c\le C<\infty$，使所有 $0<b\le b_{\rm ref}$、
$0\le\varepsilon\le1$ 及 $j\in\{\mathrm{full},\mathrm{bits},\mathrm{marked}\}$ 满足
\begin{equation}
 c\min\{b^2,\varepsilon\}\le V_j(b,\varepsilon)
 \le C\min\{b^2,\varepsilon\}.
 \label{eq:main}
\end{equation}
特别地，$V_j(b,0)=0$。常数可依赖维数、模型、$k$ 和 $b_{\rm ref}$，
但不依赖变化的 $b,\varepsilon$。
\end{theorem}

\subsection{峰值约束下的损伤上界}
令 $\mathcal P_t^-$ 为长度 $t$ 的所有事件前缀产生的先验协方差有限值集。
以下最大值均取 ${\bm{P}}\in\mathcal P_t^-$、$i\in\{0,1\}$，定义
\begin{equation}
 \begin{aligned}
 \alpha_t&=\max_{{\bm{P}},i}\|{\bm{I}}-{\bm{K}}_i({\bm{P}})\|_2,\quad
 \beta_t=\max_{\bm{P}}\|{\bm{K}}_1({\bm{P}})\|_2,\\
 r_0&=\beta_0,\quad r_t=\alpha_t\|{\bm{A}}\|_2r_{t-1}+\beta_t,\\
 \ell_0&=0,\quad \ell_t=\|{\bm{A}}\|_2r_{t-1}\quad(t\ge1).
 \end{aligned}
 \label{eq:radii}
\end{equation}
由引理\ref{lem:gaussian}递推可得 $\|{\bm{d}}_t\|_2\le\ell_tb$、$\|{\bm{D}}_t\|_2\le r_tb$。
这两个半径界对全部允许策略和事件历史统一成立。
给定事件前缀 $h$，令 ${\bm{\Sigma}}_h={\bm{P}}_t^-(h)+{\bm{R}}$、
${\bm{W}}_h=({\bm{Z}}^{-1}+{\bm{\Sigma}}_h)^{-1}$、
$s_h=\det({\bm{I}}+{\bm{Z}}^{1/2}{\bm{\Sigma}}_h{\bm{Z}}^{1/2})^{-1/2}$。高斯积分给出
\begin{equation}
 \Pr_\pi(\gamma_t=0\mid h)
 =s_h\E_\pi[e^{-{\bm{d}}_t^\top {\bm{W}}_h{\bm{d}}_t/2}\mid h].
 \label{eq:prefix-silence}
\end{equation}
无攻击时的条件静默概率为 $s_h$。由 ${\bm{W}}_h\preceq {\bm{Z}}$ 及 $1-e^{-u}\le u$，
每步概率差至多为 $\|{\bm{Z}}\|_2\ell_t^2b^2/2$。
在共同事件前缀上依次耦合两个伯努利条件分布，得到
\begin{equation}
 \TV(\Law_\pi(\gamma_{0:k}),\Law_0(\gamma_{0:k}))
 \le\frac{\|{\bm{Z}}\|_2b^2}{2}\sum_{t=1}^k\ell_t^2.
 \label{eq:prefix-tv}
\end{equation}
记 $\mathcal P_k^+$ 为后验协方差有限值集，
$\Omega_k=\max_{{\bm{P}}\in\mathcal P_k^+}\operatorname{tr}{\bm{P}}-
\min_{{\bm{P}}\in\mathcal P_k^+}\operatorname{tr}{\bm{P}}$。
对协方差迹应用有界函数期望差的TV界，再结合式\eqref{eq:decomposition}，得到
\begin{equation}
 J_k\le C_{\rm peak}b^2,\qquad
 C_{\rm peak}=r_k^2+\frac{\Omega_k\|{\bm{Z}}\|_2}{2}\sum_{t=1}^k\ell_t^2.
 \label{eq:peak-bound}
\end{equation}

\subsection{记录约束下的损伤上界}
令 ${\bm{P}}={\bm{P}}_k^+$、${\bm{D}}={\bm{D}}_k$、$Y=\ind\{\gamma_{k+1}=0\}$，定义
\begin{equation}
 \begin{aligned}
 {\bm{\Sigma}}({\bm{P}})&={\bm{A}}{\bm{P}}{\bm{A}}^\top+{\bm{Q}}+{\bm{R}},\\
 s({\bm{P}})&=\det({\bm{I}}+{\bm{Z}}^{1/2}{\bm{\Sigma}}({\bm{P}}){\bm{Z}}^{1/2})^{-1/2},\\
 w({\bm{P}})&=s({\bm{P}})^{-1},\\
 {\bm{H}}({\bm{P}})&=\tfrac12 {\bm{A}}^\top[{\bm{Z}}^{-1}+{\bm{\Sigma}}({\bm{P}})]^{-1}{\bm{A}}.
 \end{aligned}
 \label{eq:witness-functions}
\end{equation}
目标时刻的偏差 ${\bm{D}}$ 经状态转移成为下一时刻的预测偏差 ${\bm{A}}{\bm{D}}$。
给定 $\mathcal G_k^+$，下一测量相对实际预测服从 $N({\bm{A}}{\bm{D}},{\bm{\Sigma}}({\bm{P}}))$，因而
\begin{equation}
 \begin{aligned}
 \E_\pi[Y\mid\mathcal G_k^+]&=s({\bm{P}})e^{-{\bm{D}}^\top {\bm{H}}({\bm{P}}){\bm{D}}},\\
 \E_0[w({\bm{P}})Y]&=1,\quad
 \E_\pi[w({\bm{P}})Y]=\E_\pi e^{-{\bm{D}}^\top {\bm{H}}({\bm{P}}){\bm{D}}}.
 \end{aligned}
 \label{eq:witness}
\end{equation}
由式\eqref{eq:witness}可知，当前估计偏差通过反馈降低下一时刻的条件静默概率。
$s({\bm{P}})$ 为相应协方差下的正常概率，权重 $w({\bm{P}})$ 将不同事件历史的正常基准统一为1。
因此，加权静默统计量的期望差能够反映当前偏差；下一时刻的载荷修改发生在发送决定之后，不影响这一关系。
${\bm{A}}$ 可逆使全部 ${\bm{H}}({\bm{P}})\succ0$。令
\begin{equation}
 \begin{aligned}
 h_*&=\min_{{\bm{P}}\in\mathcal P_k^+}\lambda_{\min}({\bm{H}}({\bm{P}})),\quad
 M=\max_{{\bm{P}}\in\mathcal P_k^+}w({\bm{P}}),\\
 B_*&=r_kb_{\rm ref},\quad q=\frac{1-e^{-h_*B_*^2}}{B_*^2}>0.
 \end{aligned}
 \label{eq:witness-constants}
\end{equation}
凹函数 $1-e^{-h_*x}$ 在 $[0,B_*^2]$ 上的割线给出
$1-e^{-{\bm{D}}^\top {\bm{H}}({\bm{P}}){\bm{D}}}\ge q\|{\bm{D}}\|_2^2$。又 $0\le w({\bm{P}})Y\le M$，故
\begin{equation}
 \begin{aligned}
 q\E_\pi\|{\bm{D}}\|_2^2&\le\E_0[w({\bm{P}})Y]-\E_\pi[w({\bm{P}})Y]\\
 &\le M\Delta_{\rm marked}(\pi),\\
 J_k(\pi)&\le C_{\rm record}\Delta_{\rm marked}(\pi),\\
 C_{\rm record}&=M/q+\Omega_k.
 \end{aligned}
 \label{eq:record-bound}
\end{equation}
将式\eqref{eq:peak-bound}与式\eqref{eq:record-bound}结合，并利用三类记录间的数据处理关系，
即可得到定理\ref{thm:main}的上界。该上界同时计入偏差平方与协方差期望变化。

\subsection{可达到损伤下界的攻击构造}
为证明上述损伤阶数可以达到，需要构造满足因果权限与两项预算的攻击。
以下采用私有随机符号，并通过下一拍补偿控制完整日志的联合分布差异。
固定单位方向 ${\bm{v}}\in\R^n$，在目标时刻前不施加攻击。
令 $h=\gamma_{0:k-1}$（$k=0$ 时为空串），${\bm{K}}_h={\bm{K}}_1({\bm{P}}_k^-(h))$，并取
\begin{equation}
 \kappa=\bigl[\max\{1,\max_h\|{\bm{A}}{\bm{K}}_h{\bm{v}}\|_2\}\bigr]^{-1},\quad \eta=\kappa b.
 \label{eq:kappa}
\end{equation}
目标时刻发送时注入 ${\bm{a}}_k=\sigma\eta {\bm{v}}$，其中 $\sigma$ 是私有的等概率正负符号。
若下一时刻也发送，注入 ${\bm{a}}_{k+1}=\sigma\eta {\bm{A}}{\bm{K}}_h{\bm{v}}$；
目标时刻静默或其他时刻均不攻击。两个动作都满足欧氏峰值 $b$。
该策略只使用事件历史、私有符号及已知模型，因而满足因果信息约束。

记无攻击运行的前缀概率为 $w_h$、目标时刻的条件发送概率为 $p_h$。
由于目标时刻之前未施加攻击，目标时刻的协方差分布保持不变，故
\begin{equation}
 J_k=L_{\bm{v}}\eta^2,\qquad L_{\bm{v}}=\sum_hw_hp_h\|{\bm{K}}_h{\bm{v}}\|_2^2>0.
 \label{eq:lower-damage}
\end{equation}
% Essential joint-log proof; explicit constants are expanded in Appendix B.
由于两拍载荷使用同一私有符号，需要计算其联合分布，才能检验完整日志约束。
给定攻击前完整历史，令
${\bm{S}}_h={\bm{P}}_k^-(h)+{\bm{R}}$、${\bm{S}}'_h={\bm{A}}F_1({\bm{P}}_k^-(h)){\bm{A}}^\top+{\bm{Q}}+{\bm{R}}$，
${\bm{u}}_h={\bm{A}}{\bm{K}}_h{\bm{v}}$，并以 $\phi_{\bm{\Sigma}}$ 表示 $N(0,{\bm{\Sigma}})$ 的密度。
当前真实创新的未归一化发送密度为
$f_h({\bm{u}})=\phi_{{\bm{S}}_h}({\bm{u}})[1-g({\bm{u}})]$，其积分为 $p_h$。
给定当前真实发送测量，下一真实条件创新 ${\bm{e}}$ 服从 $N(0,{\bm{S}}'_h)$，
且该条件分布不随当前创新改变。
对有符号幅值 $t$，下一预测偏差为 $-t{\bm{u}}_h$；
若下一时刻发送，补偿使接收创新恰为 ${\bm{e}}$。
下一记录既可能是静默事件，也可能包含连续载荷。将两者共同表示为
$\{\bot\}\sqcup\R^n$ 上的概率核，并定义当前发送分支的联合核：
\begin{equation}
 \begin{aligned}
 \mathcal K_{h,{\bm{d}}}(\{\bot\})&=\int\phi_{{\bm{S}}'_h}({\bm{e}})g({\bm{e}}+{\bm{d}})\,d{\bm{e}},\\
 \mathcal K_{h,{\bm{d}}}(d{\bm{e}})&=\phi_{{\bm{S}}'_h}({\bm{e}})[1-g({\bm{e}}+{\bm{d}})]\,d{\bm{e}},\\
 F_h(t)(d{\bm{z}},d{\bm{\xi}})&=f_h({\bm{z}}-t{\bm{v}})\,d{\bm{z}}\,\mathcal K_{h,-t{\bm{u}}_h}(d{\bm{\xi}}).
 \end{aligned}
 \label{eq:main-joint-kernel}
\end{equation}
其中，${\bm{z}}$ 为当前接收创新，$F_h(t)$ 的总质量恒为 $p_h$。
式\eqref{eq:main-joint-kernel}的第一因子描述当前载荷的平移，第二因子描述预测反馈引起的下一事件概率变化。

记 $\|\cdot\|_{\rm var}$ 为有符号测度的变差范数，并令 $a_h={\bm{u}}_h^\top {\bm{Z}}{\bm{u}}_h$。
由 $\sup_{\bm{z}}|\partial_{{\bm{u}}_h}g({\bm{z}})|\le\sqrt{a_h/\mathrm e}$、
$\sup_{\bm{z}}|\partial_{{\bm{u}}_h}^2g({\bm{z}})|\le a_h$，
分别计入静默原子和连续密度的导数，得到
\begin{equation}
 \begin{aligned}
 \|F_h''(t)\|_{\rm var}
 &\le \|\partial_{\bm{v}}^2 f_h\|_1
 +4\sqrt{a_h/\mathrm e}\|\partial_{\bm{v}} f_h\|_1
 +2p_ha_h\\
 &\le B_h<\infty.
 \end{aligned}
 \label{eq:main-kernel-smooth}
\end{equation}
式\eqref{eq:main-kernel-smooth}中的三项依次来自当前密度的二阶导数、
两个因子的交叉导数及下一记录核的二阶导数。
$f_h$ 的高斯方向导数可积，平移保持其 $L^1$ 范数；
在共同支配测度 $d{\bm{z}}\otimes(\delta_\bot+d{\bm{\xi}})$ 下，乘积核因而二阶连续 $L^1$ 可微。
固定模型和有限事件前缀使 $B_h$ 与变化的 $b,\varepsilon$ 无关；
附录\ref{app:full-kernel}给出其显式选择。
等概率正负符号使联合核的一阶扰动相消。
对 $Q_h^\eta=[F_h(\eta)+F_h(-\eta)]/2$ 使用中心 Taylor 公式，可得
\begin{equation}
 \begin{aligned}
 &Q_h^\eta-F_h(0)\\
 &\quad=\frac12\int_0^\eta(\eta-t)
 \bigl[F_h''(t)+F_h''(-t)\bigr]\,dt,\\
 &\|Q_h^\eta-F_h(0)\|_{\rm var}\le \tfrac12B_h\eta^2.
 \end{aligned}
 \label{eq:main-kernel-taylor}
\end{equation}
给定攻击前历史及 $({\bm{z}},{\bm{\xi}})$，按
${\bm{m}}_{k+1}^-={\bm{A}}({\bm{m}}_k^-+{\bm{K}}_h{\bm{z}})$，并在 ${\bm{\xi}}\ne\bot$ 时令下一载荷为
${\bm{m}}_{k+1}^-+{\bm{\xi}}$，即可重建全部事件、载荷和预测记录。
该重建映射不依赖 $t$ 或私有符号；当前静默分支则与无攻击运行相同。
对共同前史积分，利用 TV 在固定映射下的收缩性，以及等质量测度差的
TV 为其变差范数的一半，得到
$\Delta_{\rm full}\le C_\pm\eta^2$，
其中 $C_\pm=\frac14\sum_hw_hB_h\in(0,\infty)$。

再在实验开始前以私有概率
$\lambda=\min\{1,\varepsilon/(C_\pm\kappa^2b^2)\}$ 选择上述攻击，其余情况下不施加攻击，得到
\begin{equation}
 \begin{aligned}
 \Delta_{\rm full}&\le\varepsilon,\\
 J_k&=L_{\bm{v}}\min\{\kappa^2b^2,\varepsilon/C_\pm\}.
 \end{aligned}
 \label{eq:attainable}
\end{equation}
取 $c=L_{\bm{v}}\min\{\kappa^2,C_\pm^{-1}\}$ 即得全部三类记录的下界，主定理得证。
上述构造中的下一拍动作仅补偿接收创新的平移，不消除条件估计偏差，
也不改变已经形成的目标时刻损伤。它使载荷分布与后续事件概率能够在同一联合核中分析，从而满足完整日志预算。

\begin{remark}[维数、观测条件和时间范围]
式\eqref{eq:main}对任意固定 $n$ 给出同一阶数结构，其中的常数依赖具体维数与模型。
可逆性保证下一时刻事件所对应的二次型在全部偏差方向上正定。
若 ${\bm{A}}$ 奇异，则对非零 ${\bm{v}}\in\ker {\bm{A}}$ 有 ${\bm{v}}^\top {\bm{H}}({\bm{P}}){\bm{v}}=0$；
沿该方向的当前偏差不改变下一时刻预测均值，因而不能由下一静默概率约束。
这是单步事件感知全部偏差方向所需的结构条件。
若使用部分测量 ${\bm{y}}={\bm{C}}{\bm{x}}+{\bm{v}}$，对应二次型含有 ${\bm{C}}{\bm{A}}$，
仅凭 $({\bm{A}},{\bm{C}})$ 可观不能保证单步二次型正定。
因此，本节结论适用于全状态含噪测量及固定有限时刻；
部分测量情形与时间一致界需要另行分析。
有限时刻分析允许 ${\bm{A}}$ 不稳定。
\end{remark}
